\documentclass[12pt, a4paper]{report}
\usepackage[utf8]{inputenc}
\usepackage[T1]{fontenc}
\usepackage{geometry}
\geometry{a4paper, margin=1in}
\usepackage{graphicx}
\usepackage{hyperref}
\usepackage{amsmath}
\usepackage{listings}
\usepackage{xcolor}
\usepackage{titlesec}
\usepackage{caption}
\usepackage{tabularx}
\usepackage{booktabs}

% ─────────────────────────────────────────────────────────────────────────────
%  Code-listing style. Stays readable on a monochrome print.
% ─────────────────────────────────────────────────────────────────────────────
\definecolor{codegreen}{rgb}{0,0.6,0}
\definecolor{codegray}{rgb}{0.5,0.5,0.5}
\definecolor{codepurple}{rgb}{0.58,0,0.82}
\definecolor{backcolour}{rgb}{0.95,0.95,0.92}

\lstdefinestyle{mystyle}{
    backgroundcolor=\color{backcolour},
    commentstyle=\color{codegreen},
    keywordstyle=\color{magenta},
    numberstyle=\tiny\color{codegray},
    stringstyle=\color{codepurple},
    basicstyle=\ttfamily\footnotesize,
    breakatwhitespace=false,
    breaklines=true,
    captionpos=b,
    keepspaces=true,
    numbers=left,
    numbersep=5pt,
    showspaces=false,
    showstringspaces=false,
    showtabs=false,
    tabsize=2,
    extendedchars=true,
    inputencoding=utf8,
    literate=%
      {—}{{---}}1 {–}{{--}}1
      {’}{{'}}1 {‘}{{`}}1
      {“}{{``}}1 {”}{{''}}1
      {→}{{$\to$}}1 {≈}{{$\approx$}}1
      {…}{{\ldots}}1 {×}{{$\times$}}1
}
\lstset{style=mystyle}

\lstdefinelanguage{JSON}{
    morestring=[b]",
    morestring=[d]',
    morekeywords={true,false,null},
    sensitive=true,
    morecomment=[l]{//},
    morecomment=[s]{/*}{*/},
}

% Placeholder box for figures whose assets aren't ready yet.
\newcommand{\figph}[2]{%
    \begin{figure}[h]
        \centering
        \fbox{\parbox{0.85\linewidth}{\centering \textbf{PLACEHOLDER:} #1}}
        \caption{#2}
    \end{figure}%
}

\title{End-of-Studies Project (PFE)\\[6pt] \Large A Privacy-First, No-Code AI Platform\\ \large for Tabular ML, Local Generative AI, and Deep Learning}
\author{[Student Name]\\ \small Supervised by [Supervisor Name]}
\date{\today}

\begin{document}

\maketitle

\begin{abstract}
This report describes a no-code AI platform built over six sprints as an
end-of-studies project. The platform lets non-technical users upload their
data, sketch a pipeline on a visual canvas, and train a model under one
firm constraint: nothing ever leaves the machine. Three workloads share
the same canvas. The first is classical tabular ML across fourteen
algorithms. The second is local Generative AI through a
Retrieval-Augmented Generation chat powered by Ollama and
\texttt{pgvector}. The third is deep-learning image classification on a
small CNN catalogue sized for a 6\,GB consumer GPU. The chapters that
follow walk through each sprint's design choices, the things that broke
on the first try, and the conventions I leaned on throughout: Celery for
anything CPU-bound, MongoDB for anything that resists a schema, and a
strict gateway-only header-injection scheme that gave the microservices
a zero-trust boundary without the usual gymnastics.
\end{abstract}

\tableofcontents
\newpage

\listoffigures
\newpage

% ═════════════════════════════════════════════════════════════════════════════
\chapter*{General Introduction}
\addcontentsline{toc}{chapter}{General Introduction}
% ═════════════════════════════════════════════════════════════════════════════

Machine learning is now routine in almost every industry, but the
tooling around it still defaults to one of two extremes. On one end sit
the cloud-managed services --- DataRobot, Vertex AI, Azure ML Studio ---
which hide the complexity behind polished interfaces but require users
to ship their data into someone else's infrastructure. On the other end
live the open notebooks where everything is possible and nothing is
easy: Jupyter, Colab, raw Python. The space between those two poles,
where a non-technical user might want to build a model without either
learning Python or handing over their data, is surprisingly thin.

That gap is the question this project sets out to answer. The goal was
a platform that gives a domain expert --- a clinician, an analyst, a
small business owner --- enough visual scaffolding to build, train, and
inspect a model on their own machine. Not a tutorial, not a hosted
SaaS, but a real piece of software that runs locally, accepts local
data, and never phones home to a vendor API.

Three workloads sit behind the same canvas: tabular ML, RAG-style
generative AI, and image classification. Each one imposes a different
operational profile, and fitting them into a single coherent product
forced choices that were not obvious upfront. The report is organised
as follows. Chapter~1 sets the project context. Chapter~2 reviews the
existing landscape and the problems the new platform is meant to fix.
Chapter~3 describes the methodology and the toolchain. Chapters~4
through~9 walk through the six sprints in order: the foundation sprint,
then the three AI workloads back to back (ML, LLM, DL), then the
collaboration layer, then the operability sprint that turned the
platform from ``demonstrably works'' into ``demonstrably operable''.
Chapter~10 covers testing and deployment. The General Conclusion steps
back to assess what worked, what didn't, and what I would build with
another quarter.

\newpage

% ═════════════════════════════════════════════════════════════════════════════
\chapter{General Context of the Project}
% ═════════════════════════════════════════════════════════════════════════════

\section*{Introduction}

This chapter frames the problem the platform is built to address. It
describes the host context, the platform itself in plain terms, the
functional and non-functional requirements I committed to before
writing code, the technical architecture at a high level, and the
deliverables the work was evaluated against.

\section{Project Overview and Context}

\subsection{Host context}

The project was developed as an end-of-studies (PFE) project for a
National Bachelor's Degree in Information Technologies, specialisation
Development of Information Systems. The deliverable is a working
software product plus this report. The host setting is academic rather
than corporate: the supervisor is the academic advisor, the audience
is the jury, and the demo machine is a personal laptop.

\subsection{Project description --- the no-code AI platform}

The platform is a self-hosted web application. A signed-in user can:

\begin{itemize}
    \item Upload tabular data (CSV, Excel), pull it from a Postgres or
        MySQL connector, or upload a zip of images organised by class
        folder.
    \item Inspect a profile of the dataset --- missing values, skew,
        outliers, target imbalance, correlations --- before doing
        anything else with it.
    \item Build a pipeline visually on a React Flow canvas, in one of
        three modes:
    \begin{itemize}
        \item \textbf{Traditional ML} --- fourteen supervised and
            unsupervised algorithms spanning classification,
            regression, clustering, and time-series forecasting.
        \item \textbf{Generative AI} --- index documents into a local
            vector store and chat with a RAG pipeline.
        \item \textbf{Deep Learning} --- image classification on a
            fixed catalogue of three CNN architectures.
    \end{itemize}
    \item Train, evaluate, download the trained artefact, and ---
        for image models --- run live inference on a hand-uploaded
        test image.
    \item Collaborate with teammates inside a company workspace, with
        role-based access control, real-time canvas presence, threaded
        per-pipeline chat, and one-click Google Meet links.
\end{itemize}

A super-admin role sits alongside the user-facing application with its
own dashboard for user management, audit logs, queue health, hardware
monitoring, and a ``view-as-user'' impersonation flow gated by a hard
five-minute TTL.

\figph{Landing page mockup}{Landing page interface}

\section{Requirements Specification}

\subsection{Functional Requirements}

\begin{itemize}
    \item Tabular ML pipelines: dataset $\to$ train $\to$ evaluate,
        backed by fourteen algorithms.
    \item RAG pipelines: document $\to$ vector store $\to$ chat, end
        to end on the local box.
    \item DL pipelines: image dataset $\to$ CNN architecture $\to$
        training loop, fully local.
    \item User authentication with email/password and optional
        TOTP-based two-factor authentication.
    \item Stripe-backed billing across free, solo, and company tiers,
        with super-admin overrides on per-user quotas.
    \item Real-time presence and shared editing on the pipeline
        canvas.
    \item Per-pipeline access control with five roles (Owner, Project
        Manager, Data Scientist, Analyst, Viewer).
    \item Super-admin operations: user management, audit logs, queue
        and hardware monitoring, user impersonation.
\end{itemize}

\subsection{Non-Functional Requirements}

\begin{itemize}
    \item \textbf{Data sovereignty.} No outbound calls to AI vendors.
        Stripe is the only third-party service in the deployment
        graph, and it is optional --- without an API key the platform
        still works on the free tier.
    \item \textbf{Modularity.} The system splits into five
        independent Python services so a long training run cannot
        block an authentication request.
    \item \textbf{Reproducibility.} Every dataset version is profiled,
        every training run is recorded, every model artefact is
        versioned and downloadable.
    \item \textbf{Hardware constraints.} The whole stack runs on a
        single consumer machine. My development box --- a laptop with
        $16$\,GB of RAM and a GTX~1660~Super GPU --- was picked on
        purpose, because it sits close to the floor of what someone
        might realistically deploy on.
    \item \textbf{Accessibility.} WCAG~AA contrast on the default
        theme, with a high-contrast theme available from the navbar.
    \item \textbf{Internationalisation.} English and French resource
        bundles, with TypeScript bindings generated from the English
        bundle so a missing translation fails the type-check.
\end{itemize}

\section{Technical Architecture}

\subsection{Technology Stack}

\begin{itemize}
    \item \textbf{Frontend.} React~18, Vite, TypeScript, MUI, React
        Flow, Plotly, i18next, Zustand.
    \item \textbf{Backend.} Five Flask services, Celery for
        CPU-bound work, Flask-SocketIO for live channels.
    \item \textbf{Data.} Postgres (with \texttt{pgvector}), MongoDB,
        Redis, TimescaleDB for metrics.
    \item \textbf{ML.} scikit-learn, XGBoost, LightGBM, CatBoost, H2O,
        Prophet, SHAP.
    \item \textbf{LLM.} Ollama serving Llama~3.2~3B
        (\texttt{Q4\_K\_M} quantisation), sentence-transformers
        (\texttt{all-MiniLM-L6-v2}) for embeddings.
    \item \textbf{DL.} PyTorch + torchvision, CUDA~12.1.
    \item \textbf{Ops.} Docker Compose, GitHub Actions, MailHog for
        development email capture.
\end{itemize}

\subsection{Technical Constraints}

\begin{itemize}
    \item Single host, single GPU (6\,GB VRAM).
    \item No cloud AI vendor in the request path --- a hard
        non-negotiable, not a preference.
    \item Stack must come up with a single \texttt{make up} on a
        clean machine.
    \item CI must complete in under fifteen minutes for a
        full-stack change.
\end{itemize}

\section{Deliverables and Evaluation Criteria}

\subsection{Project Deliverables}

\begin{itemize}
    \item A working web application running the three pipeline modes
        end to end on the demo machine.
    \item Five Python microservices, packaged as Docker images.
    \item A React + TypeScript frontend, packaged as a static build.
    \item A \texttt{docker-compose.yml} that brings the entire stack
        online with one command.
    \item A CI pipeline (lint + tests, parallel jobs) green on the
        \texttt{main} branch.
    \item This written report and the supporting diagrams.
\end{itemize}

\subsection{Evaluation Criteria}

\begin{itemize}
    \item Functional coverage of the three pipeline modes.
    \item Soundness of the architecture (separation of concerns,
        explicit security boundaries, reproducible deployment).
    \item Code quality (typing, linting, test coverage on the
        load-bearing paths).
    \item Quality of the report and the live demo.
\end{itemize}

\section{Interface Mockups}

\figph{Login + 2FA screen}{Authentication interface (email/password + optional TOTP)}

\figph{Home / Dashboard screen}{Home page interface for a signed-in user}

\figph{Admin operations dashboard}{Admin dashboard interface}

\section{Project Duration}

The project ran across six sprints of roughly two to three weeks each,
with a runnable demo and a self-contained git tag at the end of every
sprint. The timeline, with start and end dates, is captured in the
Gantt chart in Chapter~3.

\section*{Conclusion}

The chapter framed the problem the rest of the report addresses --- how
to make ML approachable without sacrificing data sovereignty --- and
sketched the platform that emerged in response. The next chapter
reviews the existing landscape to explain why the question is harder to
answer with off-the-shelf tools than one might initially assume.

% ═════════════════════════════════════════════════════════════════════════════
\chapter{Preliminary Study}
% ═════════════════════════════════════════════════════════════════════════════

\section*{Introduction}

Before describing how the platform was built, it helps to look at why
nothing already on the market quite fits the problem. The no-code ML
space is crowded, but the products in it cluster into families that
all share the same blind spot.

\section{Current Situation}

Three categories of product dominate the space:

\begin{itemize}
    \item \textbf{Cloud-native AutoML platforms} --- DataRobot, H2O
        Driverless~AI, Azure ML Studio, Google Vertex~AI. Highly
        polished, exhaustive algorithm catalogues, almost
        invariably SaaS-only. Their privacy story ends at ``trust
        our infrastructure''. For some buyers that is fine; for
        regulated industries it isn't.
    \item \textbf{Open-source notebooks} --- Jupyter, Google Colab,
        DataLab. Powerful for experts, but every workflow
        eventually bottoms out in raw Python. That is exactly the
        friction the no-code category was invented to remove.
    \item \textbf{Visual ETL tools} --- KNIME, Alteryx, RapidMiner.
        Closer in spirit to this project, but their centre of
        gravity is data preparation rather than model training,
        and their LLM integrations tend to route requests back
        through OpenAI.
\end{itemize}

\figph{Comparison table of existing solutions}{Positioning of existing no-code ML solutions}

\section{Observed Problems}

\subsection{Forced data exfiltration}

A clinician analysing patient records or a bank analyst classifying
transactions cannot lawfully upload that data to a vendor bucket. The
problem isn't contractual --- it is regulatory (GDPR, HIPAA, local
banking rules). Every SaaS AutoML platform in the previous section
assumes the user is fine with that upload. They are not.

\subsection{Cloud LLM dependence in generative-AI tooling}

Visual tools that have added an ``AI'' panel in the last two years
almost always route the request through OpenAI or a similar API. The
panel is bolted on, not local. For a buyer who can't ship the prompt
out either, the panel is unusable.

\subsection{The notebook cliff}

Notebooks solve the locality problem cleanly. They fail on the
no-code one. A domain expert who can read a confusion matrix but
cannot write \texttt{from sklearn.ensemble import RandomForestClassifier}
hits a wall on line one of the first cell. Closing that gap is the
whole reason a visual canvas exists in the first place.

\subsection{Hardware ambiguity}

The few self-hostable products that do exist (KubeFlow, MLflow on
local Kubernetes) assume the operator has a Kubernetes cluster
sitting around. That is not the operating environment of a small
team that has a workstation and an ambition.

\section{Objectives of the New System}

The platform I built sets out to satisfy four objectives at once:

\begin{enumerate}
    \item \textbf{Locality.} Every model server, vector store, queue,
        and database runs inside the same Docker Compose stack on the
        user's own machine.
    \item \textbf{Coverage.} Three workloads --- tabular ML, RAG-style
        generative AI, and image classification --- share the same
        canvas. The user does not need a different product for each.
    \item \textbf{Approachability.} A canvas-based UI that a
        non-programmer can drive end to end, with sensible
        per-node validation and informative error messages on the
        unhappy path.
    \item \textbf{Honest hardware story.} The reference deployment is
        a laptop with $16$\,GB of RAM and a GTX~1660~Super (6\,GB
        VRAM). The platform must run on it, not just compile on it.
\end{enumerate}

Two technical shifts make the fourth objective realistic today and
would not have made it realistic two years ago.

\begin{enumerate}
    \item \textbf{Aggressive LLM quantisation.} Four-bit quantisation
        (\texttt{Q4\_K\_M}) shrinks a three-billion-parameter model
        from roughly $12$\,GB of weights down to about $2$\,GB,
        while preserving most of the original quality on
        conversational tasks. The trade-off is real but, for chat,
        it isn't fatal.
    \item \textbf{First-class GPU inference runtimes for consumer
        hardware} (Ollama, llama.cpp, vLLM). A 1660 Super, with
        $6$\,GB of VRAM and $30$\,TFLOPS of FP16 throughput, can
        serve roughly $30$ tokens per second on a 3B model ---
        enough latency budget for an interactive RAG chat.
\end{enumerate}

\figph{Quantisation $\to$ VRAM mapping}{Mapping of quantisation level to VRAM footprint for a 3B-parameter model}

\section*{Conclusion}

Cloud platforms solve the wrong half of the problem, notebooks demand
too much of the user, and visual ETL tools are not aimed at model
training. The platform described in the rest of this report fills the
gap deliberately: it is local by construction, broad enough to cover
the three AI workloads that matter today, and modest enough in its
hardware appetite to actually run on a laptop.

% ═════════════════════════════════════════════════════════════════════════════
\chapter{Methodology and Working Environment}
% ═════════════════════════════════════════════════════════════════════════════

\section*{Introduction}

Six sprints, one person, three AI workloads on a single GPU --- the
project's main risk was not technical. It was scope management. This
chapter explains the methodology I used to keep that risk in check,
the toolchain that anchored the build, and the layout of the services
those choices produced.

\section{Comparison of Project Management Methodologies}

I considered three styles before settling.

\begin{itemize}
    \item \textbf{Waterfall.} Plan everything upfront, build everything
        in order, demo at the end. Wrong shape for an ML project
        where the cost of a feature is genuinely unpredictable
        (Prophet's CmdStan dependency added two hours to a
        ``trivial'' algorithm).
    \item \textbf{Pure Kanban.} Continuous flow, no sprint boundary.
        Tempting for a solo project, but it removes the
        checkpoints I needed to demo to my supervisor.
    \item \textbf{Scrum-lite (chosen).} Short sprints, a runnable
        demo at the end of each, a self-contained git tag, a
        retrospective note for myself. Enough ceremony to enforce
        scope discipline; not so much that ceremony eats the
        sprint.
\end{itemize}

\section{The Choice of Agile Methodology}

Two reasons pushed me toward Agile. First, ML capacity planning is
unpredictable. A feature that looks easy on paper can hide a
compile-from-source dependency, and Agile's sprint boundary absorbs
that kind of slip. Second, I needed dated checkpoints I could demo
without hand-waving. Without sprint structure I would have ended up
with a single ``everything works at once'' deliverable, no story to
tell along the way, and no fallback position if a late feature fell
through.

\section{Scrum Methodology}

\figph{Scrum framework diagram}{The Scrum framework: roles, artefacts, events}

\subsection{Scrum Roles}

In a one-person team the roles all collapse onto the same person.
What does not collapse is the \emph{discipline} of switching hats.
When planning, I am the Product Owner --- I write user stories and
prioritise. When building, I am the Developer. When reviewing the
sprint, I am the Scrum Master --- what went wrong, what would I do
differently. Treating those as three roles, even when they live in
one head, kept me honest.

\subsection{Scrum Artefacts}

\begin{itemize}
    \item \textbf{Product Backlog.} A single Markdown file under
        \texttt{docs/backlog.md}, grouped by sprint.
    \item \textbf{Sprint Backlog.} The slice of the product backlog
        promised for the current sprint, lifted into the sprint's
        section.
    \item \textbf{Increment.} The git tag at the end of the sprint,
        plus the demo recording.
\end{itemize}

\subsection{Scrum Events}

\begin{itemize}
    \item \textbf{Sprint Planning.} Half a day at the start of each
        sprint. Carry-over from the previous sprint plus a fresh
        slice of the backlog.
    \item \textbf{Daily check-in.} A short written log entry, not a
        meeting.
    \item \textbf{Sprint Review.} A live demo with my supervisor at
        sprint end.
    \item \textbf{Sprint Retrospective.} A retro note in the same
        Markdown backlog, under the sprint heading.
\end{itemize}

\section{Product Backlog}

The six sprints below each ended with a runnable demo and a
self-contained git tag.

\begin{itemize}
    \item \textbf{Sprint~1 --- Authentication, Gateway, Data
        Ingestion.} Foundational service mesh.
    \item \textbf{Sprint~2 --- Visual Pipeline Editor and Tabular
        ML.} Canvas, fourteen algorithms, SHAP, model export.
    \item \textbf{Sprint~3 --- Local Generative AI and RAG.} Ollama,
        \texttt{pgvector}, streamed chat.
    \item \textbf{Sprint~4 --- Deep Learning for Image
        Classification.} PyTorch service, three CNNs, the VRAM
        guard.
    \item \textbf{Sprint~5 --- Multi-tenancy and Real-Time
        Collaboration.} Project ACL, cursors, threaded chat, Meet
        links.
    \item \textbf{Sprint~6 --- Dashboard, i18n, Accessibility, Admin
        Console.} Operability and polish.
\end{itemize}

\section{Project Planning}

\figph{Gantt chart}{Gantt chart for the six-sprint plan}

The Gantt chart in the figure above tracks planned versus actual
duration. Two sprints slipped by roughly a week each --- Sprint~2
because the live training visualisation had a Plotly autoscale bug
that took a day to find, and Sprint~4 because the static VRAM
estimate needed an empirical recalibration pass on real hardware.

\section{Technologies}

\begin{itemize}
    \item \textbf{Frontend.} React~18 + Vite + TypeScript + MUI. React
        for the ecosystem, Vite because cold starts under
        \texttt{create-react-app} were embarrassingly slow, MUI
        because it ships accessible components out of the box ---
        a point that matters for the company-tier sale.
    \item \textbf{Backend.} Flask + Celery + Flask-SocketIO. Flask
        was the lowest-common-denominator choice for a
        microservice. FastAPI would have given me typed routes
        for free, but every template I started from assumed
        Flask, and going back through five rewrites just for
        typing wasn't a fight I wanted to have. Celery handles
        every CPU-bound task: profiling, preprocessing, H2O
        training, RAG ingestion, image extraction, and PyTorch
        training each live on their own queue.
    \item \textbf{Postgres + pgvector for relational and vector
        data.} Co-locating the vector store with the auth schema
        let me avoid spinning up a sixth backing service.
        \texttt{pgvector}'s IVFFlat index is more than enough for
        the few hundred thousand chunks that a single-tenant
        local deployment is ever likely to hold.
    \item \textbf{MongoDB for everything schemaless.} Datasets,
        pipelines, model versions, task results --- anything
        whose shape evolves between sprints lives in Mongo.
    \item \textbf{Redis for queue + cache + socket message bus.}
        Three logical databases on the same instance: DB~0 for
        Celery, DB~1 for Celery results, DB~2 for the SocketIO
        message queue.
    \item \textbf{Ollama for local LLM inference.} The alternative
        was calling \texttt{llama.cpp} directly. That would have
        given me finer control at the cost of building my own
        HTTP shim around it; Ollama's REST API is good enough
        that the trade wasn't close.
    \item \textbf{PyTorch + torchvision for image classification.}
        The CUDA wheels are downloaded from PyTorch's CUDA~12.1
        index inside the Dockerfile so the host only needs the
        NVIDIA driver plus \texttt{nvidia-container-toolkit}.
\end{itemize}

\section{Architecture}

\figph{Microservice topology diagram}{High-level microservice topology}

\begin{itemize}
    \item \texttt{api-gateway} (port 8000) --- decodes JWTs, injects
        \texttt{X-User-Id} / \texttt{X-User-Tier} /
        \texttt{X-User-Role} headers, proxies to the upstream
        services. Also hosts the SocketIO server for chat and
        training presence.
    \item \texttt{auth-service} (port 8001) --- user accounts,
        billing, admin operations, project and company ACL.
    \item \texttt{data-ingestion-service} (port 8002) --- file
        uploads, SQL and S3 connectors, profiling, preprocessing,
        RAG ingestion, image-dataset extraction.
    \item \texttt{ml-training-service} (port 8003) --- the fourteen
        tabular algorithms, model registry, RAG chat orchestration.
    \item \texttt{metrics-service} (port 8004) --- TimescaleDB-backed
        hardware-telemetry collection.
    \item \texttt{dl-training-service} (port 8005) --- PyTorch image
        classification.
\end{itemize}

A GitHub Actions workflow runs four jobs in parallel on every push:
backend lint (\texttt{black}, \texttt{isort}, \texttt{flake8}),
frontend lint (\texttt{eslint}, \texttt{tsc --noEmit}), backend tests
(a Postgres + Mongo + Redis stack is spun up as services, then
\texttt{pytest} runs against each microservice), and frontend tests
(\texttt{vitest} with coverage upload). End-to-end CI runs in roughly
six minutes for a frontend-only change, nine minutes for a backend
change, and fourteen minutes when both sides move at once.

\section*{Conclusion}

The methodology and toolchain held throughout the six sprints. Some
choices --- Flask in particular --- are ones I'd revisit if starting
over, but they got me to a working system without rework. The next
chapter starts the sprint walkthrough.

% ═════════════════════════════════════════════════════════════════════════════
\chapter{Sprint 1: Authentication, Gateway, and Data Ingestion}
% ═════════════════════════════════════════════════════════════════════════════

\section*{Introduction}

Sprint~1 was the foundation sprint. Nothing glamorous --- no canvases,
no models --- but the design choices made in these two weeks shaped
every later sprint. The gateway pattern and the shared-volume
convention introduced here are the quiet load-bearing decisions of
the whole project.

\section{Sprint Backlog}

\begin{itemize}
    \item Stateless JWT authentication with refresh-token rotation.
    \item Redis-backed token revocation list.
    \item Role-based access control (super-admin, tenant roles).
    \item API gateway that strips client-supplied identity headers
        and injects authoritative ones from the decoded JWT.
    \item Optional TOTP two-factor authentication.
    \item Per-IP and per-user rate limiting at the gateway.
    \item \texttt{data-ingestion-service}: file uploads, SQL
        connectors, asynchronous profiling.
    \item Mongo-backed dataset metadata.
    \item Parquet-based preprocessing output (train/val/test split).
\end{itemize}

\section{Design}

\figph{Use case diagram for Sprint 1}{Use case diagram --- authentication and ingestion}

\figph{Login + JWT issuance sequence}{Sequence diagram --- login + refresh-token rotation}

\figph{Asynchronous upload sequence}{Sequence diagram --- upload + asynchronous profiling}

\figph{Class diagram for Sprint 1}{Class diagram --- auth-service + data-ingestion-service}

\subsection{Why an opaque-but-decoded gateway}

Two reasonable architectures were on the table:

\begin{enumerate}
    \item Have the gateway hit \texttt{auth-service} on every
        request to validate the JWT.
    \item Have the gateway hold the same \texttt{JWT\_SECRET\_KEY}
        as \texttt{auth-service} and decode the token in-process.
\end{enumerate}

I went with the second, and I'd make the same call again. The gateway
is on the hot path; a synchronous HTTP round-trip per request adds
latency that compounds badly under load. The price is that the
gateway needs the same signing secret as the auth service --- which,
in a shared environment-file deployment, is a non-issue. In a
horizontally scaled production setup it would be more involved; for
the demo target it's the right trade.

\subsection{Header-injection rules}

The gateway strips any client-supplied \texttt{X-User-*},
\texttt{X-Company-*}, or \texttt{X-Project-*} header before
forwarding, then re-adds them from the decoded JWT. This is, without
exaggeration, the most important single security invariant in the
whole system. \textbf{Downstream services trust \texttt{X-User-Id}
unconditionally, which means an attacker who can supply that header
without going through the gateway can impersonate anyone.} The
strip-then-inject pattern is deliberate, well tested, and not
negotiable.

\subsection{The shared-volume convention}

My first instinct on the ingestion side was to have the API stream
files to the worker over the network. I started building that, then
realised halfway in that my Docker Compose stack already had a named
volume (\texttt{uploaded\_files}) that both the API container and the
worker container could mount. Mounting it in both places collapsed
the streaming-upload protocol into a one-line write: the API saves
to \texttt{/uploads/<id>.csv}, the worker reads from the same path.
That pattern then became the convention for every binary artefact in
the system --- trained models, RAG documents, image-dataset extracts.

\subsection{Why MongoDB for dataset metadata}

A dataset's profile is fundamentally schemaless. A CSV with two
numeric columns produces a profile that looks nothing like the
profile of a CSV with thirty mixed-type columns including a target.
I spent the better part of a day considering a Postgres JSONB column
for this and ultimately preferred Mongo's storage model. It's
lossless, it's queryable enough for the use case, and it doesn't
impose a schema migration every time I want to add a new statistic.

\section{Implementation}

\begin{lstlisting}[language=Python, caption=Gateway proxy with header injection (\texttt{api-gateway/app/routes/proxy.py})]
def _forward(upstream_base, path, require_auth=True, timeout=30):
    if require_auth:
        try:
            verify_jwt_in_request()
            claims = get_jwt()
            g.user_id = get_jwt_identity()
            g.user_role = claims.get("role", "")
            g.user_tier = claims.get("tier", "free")
        except Exception as e:
            return jsonify({"error": "unauthorized", "message": str(e)}), 401

    # Strip every client-supplied identity header before forwarding.
    _STRIP_PREFIXES = ("x-user-", "x-company-", "x-project-")
    headers = {
        k: v for k, v in request.headers
        if k.lower() not in _HOP_BY_HOP
        and k.lower() != "host"
        and not k.lower().startswith(_STRIP_PREFIXES)
    }
    # Re-inject the authoritative ones from the decoded JWT.
    if require_auth or hasattr(g, "user_id"):
        headers.update(get_forwarded_headers())

    resp = requests.request(
        method=request.method, url=f"{upstream_base.rstrip('/')}{path}",
        headers=headers, data=request.get_data(), stream=True, timeout=timeout,
    )
    return _stream_response(resp)
\end{lstlisting}

\begin{lstlisting}[language=Python, caption=Profiling task (\texttt{data-ingestion-service/app/tasks/profiling.py})]
@celery.task(name="app.tasks.profiling.profile_dataset", bind=True)
def profile_dataset(self, dataset_id: str):
    client = MongoClient(os.environ["MONGO_URL"])
    db = client[os.environ.get("MONGO_DB", "nocode_ingestion")]
    datasets = db["datasets"]
    task_results = db["task_results"]

    task_results.update_one(
        {"task_id": self.request.id},
        {"$set": {"status": "running", "progress_pct": 0}},
        upsert=True,
    )
    doc = datasets.find_one({"dataset_id": dataset_id})
    df = load_dataframe(doc["file_path"])

    summary = compute_profile_summary(df, target_column=doc.get("target_column"))

    datasets.update_one(
        {"dataset_id": dataset_id},
        {"$set": {
            "status": "ready",
            "profiling_summary": summary,
            "row_count": len(df),
            "column_count": len(df.columns),
        }},
    )
    task_results.update_one(
        {"task_id": self.request.id},
        {"$set": {"status": "success", "progress_pct": 100}},
    )
\end{lstlisting}

\subsection{Security posture}

\begin{itemize}
    \item Access tokens have a fifteen-minute TTL. Refresh tokens
        live as HttpOnly cookies and rotate on every refresh; the
        previous token's JTI lands in the Redis blacklist so a
        stolen refresh token can actually be revoked.
    \item Passwords use \texttt{bcrypt} at cost~12. Refresh tokens
        themselves are stored as SHA-256 hashes in Postgres --- the
        database-breach scenario yields a hash, not a usable
        token.
    \item TOTP-based two-factor authentication is optional per user;
        super-admins can require it per-tier from the admin panel.
    \item Per-IP and per-user rate limits sit at the gateway. The
        per-user limit defaults to $600$\,req/min in development,
        deliberately loose because React 18 Strict Mode
        double-invokes \texttt{useEffect} and a tighter cap was
        producing spurious 429s.
\end{itemize}

\subsection{SQL connectors and Fernet encryption}

Users can configure a Postgres or MySQL connector inside the
platform; the credentials are encrypted with a Fernet key
(AES-128-CBC plus HMAC) held in the service's environment, then
stored in Mongo. The same key is reused for any other secret the
ingestion service needs to hold, and rotating it is a documented
operations procedure rather than a fire-fight. SQL queries themselves
run through SQLAlchemy with parameter binding, which makes SQL
injection structurally impossible rather than merely unlikely.

\figph{Login screen}{Login screen with optional TOTP field}

\figph{Connector wizard}{SQL connector wizard --- credentials entered, encrypted at rest}

\section{Unit Tests}

The auth-service test suite sits above $80\%$ line coverage on the
authentication path and the admin operations. Notable cases:

\begin{itemize}
    \item Forged \texttt{X-User-Id} header: the request goes through
        the gateway, the forged header is stripped, the
        authoritative one from the JWT replaces it. Asserted on a
        synthetic upstream that echoes back the headers it saw.
    \item Refresh-token replay: a refresh token is used once, then
        replayed. The second use is rejected by the blacklist.
    \item Expired access token: the gateway returns 401 with the
        documented error envelope.
    \item Rate limit: 601 requests in one minute, the 601st returns
        429 with a \texttt{Retry-After} header.
\end{itemize}

The ingestion suite uses \texttt{mongomock} so it does not need a
real Mongo to run. It exercises the upload-then-profile path
end-to-end with a small fixture CSV and asserts the profile shape.

\section*{Conclusion}

By the end of Sprint~1 the platform had end-to-end authentication, a
working refresh-token loop, a gateway pattern the next five sprints
would build on, and an ingestion service that could profile
arbitrarily messy CSVs asynchronously. The
authoritative-headers convention turned out to be one of the
project's most quietly useful design decisions: every subsequent
service trusts \texttt{X-User-Id} without a second thought, because
the gateway is the only path along which it can be set.

% ═════════════════════════════════════════════════════════════════════════════
\chapter{Sprint 2: Visual Pipeline Editor and Tabular Machine Learning}
% ═════════════════════════════════════════════════════════════════════════════

\section*{Introduction}

Sprint~2 was the largest single sprint of the project. It pulled
together four big pieces --- the visual pipeline editor, the
\texttt{ml-training-service} with its fourteen tabular algorithms,
the live metric visualisations, and the post-training surface (SHAP
explainability, model export) --- into a coherent flow the user
could actually follow from canvas to model.

\section{Sprint Backlog}

\begin{itemize}
    \item React Flow canvas with custom node components.
    \item Pipeline schema (Mongo document with nodes + edges).
    \item \texttt{ml-training-service} hosting fourteen tabular
        algorithms.
    \item \texttt{BaseMLModel} abstraction and an algorithm-agnostic
        Celery training task.
    \item Model registry with versioning.
    \item Live metric streaming over SocketIO.
    \item SHAP explainability for supervised runs.
    \item Confusion matrices, residual plots.
    \item Model export (joblib + generated FastAPI inference script).
\end{itemize}

\section{Design}

\figph{Use case diagram for Sprint 2}{Use case diagram --- canvas + tabular ML}

\figph{Class diagram for Sprint 2}{Class diagram --- BaseMLModel hierarchy + model registry}

\figph{Pipeline editor screenshot --- ML mode}{Pipeline editor in tabular ML mode}

\subsection{Graph schema}

A pipeline is stored in Mongo as a document with two arrays:

\begin{lstlisting}[language=JSON, caption=Pipeline document shape]
{
  "pipeline_id": "...",
  "user_id": "...",
  "type": "ml" | "rag" | "dl",
  "nodes": [
    { "node_id": "n1", "type": "dataset", "position": {"x": 60, "y": 120},
      "data": { "dataset_id": "..." } },
    { "node_id": "n2", "type": "train", "position": {"x": 320, "y": 120},
      "data": { "algorithm": "xgboost", "task_type": "classification",
                "hyperparams": {...}, "target_column": "..." } },
    { "node_id": "n3", "type": "evaluate", "position": {"x": 580, "y": 120},
      "data": {} }
  ],
  "edges": [
    { "source": "n1", "target": "n2" },
    { "source": "n2", "target": "n3" }
  ],
  "status": "draft"|"running"|"done"|"error"
}
\end{lstlisting}

\subsection{The \texttt{BaseMLModel} abstraction}

Fourteen algorithms with one orchestrator implies an abstraction.
Rather than over-engineer it, I went with a deliberately small
abstract base class: each algorithm subclass exposes \texttt{train},
\texttt{predict}, and \texttt{evaluate} with a shared signature,
whether it's wrapping scikit-learn, XGBoost, LightGBM, CatBoost, or
Prophet. The training Celery task is algorithm-agnostic --- it
constructs a model from a registry, trains it, asks it for metrics,
and serialises the estimator. Adding a new algorithm is, in
practice, a one-file change.

\begin{lstlisting}[language=Python, caption=Algorithm-agnostic training core]
class BaseMLModel(ABC):
    def __init__(self, hyperparams: dict[str, Any]):
        self.hyperparams = hyperparams
        self._estimator = None

    @abstractmethod
    def train(self, X_train, y_train=None) -> None: ...
    @abstractmethod
    def predict(self, X) -> Any: ...
    @abstractmethod
    def evaluate(self, X_test, y_test=None) -> dict[str, Any]: ...

    @property
    def estimator(self):
        return self._estimator
\end{lstlisting}

The registry mapping algorithm strings to subclasses lives in
\texttt{ml-training-service/app/services/model\_factory.py}.

\subsection{H2O for the heavy algorithms}

For XGBoost, GBM, and Random Forest, I used H2O-3 instead of the
scikit-learn implementations. H2O's algorithms run faster on tabular
data, and they expose feature importance and SHAP values out of the
box --- which, on its own, removed an entire subtask from the
explainability work. Not every dependency pays for itself; this one
did.

\section{Implementation}

\subsection{Frontend canvas}

\figph{Node component anatomy}{Anatomy of a canvas node component --- header, body, ports, validation badge}

The canvas is a React Flow instance with custom node components.
Each node type validates itself against the surrounding graph: the
\texttt{Train} node shows an error badge if no \texttt{Dataset} node
is upstream; the \texttt{Evaluate} node shows a warning if no
training run has produced a version yet. The validation function is
pure --- \texttt{(node, edges, allNodes) $\to$ \{status, message\}}
--- which keeps it trivially testable.

\subsection{Live training visualisation}

Training emits per-step metric points over the SocketIO message
queue. The \texttt{LiveTrainingChart} component subscribes to the
pipeline's room (\texttt{pipeline\_<id>}) on the \texttt{/training}
namespace and plots the points on a Plotly scatter. The Y-axis is
auto-scaled because progress percentage ($0$--$100$) and metric
values (typically $0$--$1$) live on different ranges --- progress on
the primary axis, metrics on the secondary. A small thing, but the
chart was unreadable until I got that right.

\figph{Live training chart}{Live training chart --- progress on the primary axis, metric on the secondary}

\subsection{SHAP and the explainability surface}

Every supervised training run computes SHAP values over a sampled
subset of the test set. The bar chart of mean absolute SHAP values
surfaces in the \texttt{Evaluate} node's results panel. For
tree-based models --- XGBoost, Random Forest, GBM, LightGBM,
CatBoost --- the SHAP computation uses the optimised
\texttt{TreeExplainer} path, which is roughly two orders of
magnitude faster than \texttt{KernelExplainer}. On a $5{,}000$-row
test set, that's the difference between SHAP being a feature and
SHAP being a wait.

\figph{SHAP feature importance chart}{SHAP feature importance bar chart}

\subsection{Model export}

The export step packages the trained \texttt{.joblib} artefact along
with a generated FastAPI inference script. The user downloads a zip
they can drop onto any container runtime to serve their model
without the platform at all. That property --- being able to leave
--- mattered to me on principle. A no-code platform that locks you
in isn't really removing friction; it's relocating it.

\subsection{Why \texttt{joblib} and not \texttt{pickle}}

\texttt{joblib} was the deliberate choice for the tabular path
because its custom pickler is optimised for objects containing large
NumPy arrays --- which every scikit-learn estimator does --- and can
memory-map them on load. For a Random Forest of $500$ trees the size
difference was roughly $8\times$ in my measurements; the load-time
difference was larger. For the deep-learning path in Sprint~4 I
used \texttt{torch.save} on a \texttt{state\_dict()} instead.
Different ecosystem, different convention.

\section{Unit Tests}

\begin{itemize}
    \item Each \texttt{BaseMLModel} subclass has a happy-path test on
        a tiny fixture dataset (50 rows, four columns, a binary
        target). Assertions on output shape and on the presence of
        the documented metrics keys.
    \item The algorithm-factory is exercised by parametrised tests
        --- one per algorithm string --- that construct the model,
        train, predict, and inspect the result envelope.
    \item Node validation: parametrised tests over
        (node-type, upstream-edges, expected-status) tuples.
    \item Export bundle: the generated zip is unzipped to a
        \texttt{tempdir} and the FastAPI script is imported to
        confirm it parses. The actual inference call is exercised
        by an end-to-end test in the deployment chapter.
\end{itemize}

\section*{Conclusion}

Sprint~2 was where the project first felt like a real product. A
user could open the canvas, drop three nodes, hit Run, and watch a
model train in front of them. The \texttt{BaseMLModel} abstraction
introduced here later absorbed the deep-learning extension in
Sprint~4 without disturbing the orchestration code, and the SHAP
+ export surface closed the loop between training a model and being
able to do something with it.

% ═════════════════════════════════════════════════════════════════════════════
\chapter{Sprint 3: Local Generative AI and Retrieval-Augmented Chat}
% ═════════════════════════════════════════════════════════════════════════════

\section*{Introduction}

This was the most technically ambitious sprint of the project. A
fully local Retrieval-Augmented Generation pipeline meant fitting an
embedding model, a vector store, and an LLM inside the same Docker
stack as the rest of the platform --- and making the chat experience
responsive enough that users wouldn't notice they were running it on
a consumer GPU.

\section{Sprint Backlog}

\begin{itemize}
    \item Document ingestion pipeline (PDF / DOCX / text).
    \item Local embedding via \texttt{sentence-transformers}.
    \item Vector store in \texttt{pgvector} (IVFFlat index).
    \item Ollama integration serving Llama~3.2~3B
        (\texttt{Q4\_K\_M}).
    \item Chat orchestration: embed question, retrieve top-$k$
        chunks, build prompt, stream the answer back.
    \item Streamed UI: tokens render as they arrive over
        SocketIO.
    \item RAG-mode toggle on the canvas with three dedicated node
        types (\texttt{DocumentNode}, \texttt{VectorStoreNode},
        \texttt{ChatNode}).
\end{itemize}

\section{Design}

\figph{Use case diagram for Sprint 3}{Use case diagram --- RAG pipeline}

\figph{RAG pipeline sequence diagram}{Sequence diagram --- document $\to$ embed $\to$ store $\to$ chat}

\figph{Class diagram for Sprint 3}{Class diagram --- RAG ingestion + chat orchestrator}

The component picks:

\begin{itemize}
    \item \textbf{\texttt{sentence-transformers/all-MiniLM-L6-v2}}
        for embeddings --- 384 dimensions, around $22$\,MB on disk,
        fast enough on CPU that I never had to think about GPU
        scheduling for it.
    \item \textbf{\texttt{pgvector}} for the vector store, with an
        IVFFlat index. Co-locating it in the auth-service Postgres
        instance avoided a sixth backing service.
    \item \textbf{Ollama running \texttt{llama3.2:3b}} as the LLM
        engine. The \texttt{Q4\_K\_M} quantisation fits in roughly
        $2$\,GB of VRAM, which leaves headroom on the 6\,GB demo
        GPU for embedding, canvas, and the training queue to share
        the same device without stepping on each other.
\end{itemize}

\section{Implementation}

\begin{lstlisting}[language=Python, caption=Document ingestion task (\texttt{ml-training-service/app/tasks/rag\_ingest.py})]
@celery.task(name="app.tasks.rag_ingest.process_rag_document", bind=True)
def process_rag_document(self, document_id, pipeline_id, file_path):
    from langchain_text_splitters import RecursiveCharacterTextSplitter
    from langchain_community.document_loaders import PyPDFLoader

    docs = PyPDFLoader(file_path).load()
    splitter = RecursiveCharacterTextSplitter(
        chunk_size=500, chunk_overlap=50, length_function=len,
    )
    split_docs = splitter.split_documents(docs)
    chunk_texts = [d.page_content for d in split_docs if d.page_content.strip()]

    # Local embedding via sentence-transformers — no network call.
    vectors = embed_texts(chunk_texts)
    insert_chunks(
        pipeline_id=pipeline_id,
        document_id=document_id,
        chunks=chunk_texts,
        embeddings=vectors,
    )
\end{lstlisting}

\subsection{Streamed chat}

The chat response streams token by token over SocketIO. The
\texttt{ml-training-service} plays producer: it embeds the question,
queries pgvector, builds the prompt, and forwards Ollama's
\texttt{stream=true} chunks straight into the
\texttt{/chat/<pipeline\_id>} room. The frontend appends them to the
current assistant message as they arrive. Watching that for the
first time, on a laptop, with no network involved, was one of the
more satisfying moments of the project.

\figph{RAG chat screenshot}{RAG chat panel --- tokens stream in as the model emits them}

\subsection{Chunk sizing trade-offs}

Five hundred characters with a fifty-character overlap is the
default chunking. I tried shorter chunks (256, 128) and longer ones
(1024). The short variants improved retrieval precision on
short-answer questions and ruined it for anything that needed
multi-sentence context; the long variants did the opposite. Five
hundred sat in the middle. It is a defensible default, not a
universal answer; the project backlog has a note to expose chunk
size as a per-pipeline setting.

\section{Unit Tests}

\begin{itemize}
    \item Chunking determinism --- the same input file produces the
        same chunks across runs.
    \item Embedding shape --- 384 dimensions, normalised.
    \item Retrieval top-$k$ on a deterministic five-chunk corpus
        with hand-labelled relevance.
    \item Chat orchestrator: Ollama is mocked at the HTTP boundary,
        the orchestrator is asserted to emit the right SocketIO
        events in the right order.
\end{itemize}

The Ollama process itself is not exercised in unit tests --- it is
covered by the smoke-test script in the deployment chapter.

\section*{Conclusion}

Sprint~3 proved the local-generative-AI thesis to my own
satisfaction. The chat is not as fast as GPT-4 over the wire, and it
never will be on this hardware, but it's interactive enough to be
useful and it keeps the user's documents on their machine. For the
regulated industries the platform was designed for, that trade-off
is exactly the one they're looking to make.

% ═════════════════════════════════════════════════════════════════════════════
\chapter{Sprint 4: Deep Learning for Image Classification}
% ═════════════════════════════════════════════════════════════════════════════

\section*{Introduction}

This sprint added the third pipeline mode. Deep learning on the
demo hardware is a constrained problem --- 6\,GB of VRAM does not
forgive carelessness --- so the design started from the constraint
and worked backwards. The result is a small, opinionated catalogue
rather than a kitchen sink, and that opinionatedness is the feature.

\section{Sprint Backlog}

\begin{itemize}
    \item New microservice: \texttt{dl-training-service}.
    \item Three CNN architectures (LeNet, ResNet-9, MobileNet-V3-Small).
    \item Image-dataset ingestion (zip of \texttt{<class>/<file>}).
    \item Static VRAM estimator (the ``VRAM guard'').
    \item Tier-aware ceilings on epochs and batch size.
    \item Live training visualisation (reused from Sprint~2).
    \item Inline single-image inference (the ``Try it'' panel).
    \item DL-mode canvas toggle with three new node components.
\end{itemize}

\section{Design}

\figph{Use case diagram for Sprint 4}{Use case diagram --- DL pipeline}

\figph{DL service architecture}{Architecture --- \texttt{dl-training-service}, Celery worker, GPU allocation}

\figph{Class diagram for Sprint 4}{Class diagram --- architecture catalogue + VRAM guard}

\subsection{Scope decisions}

What's in:

\begin{itemize}
    \item Three architectures: LeNet (smoke-test scale), a custom
        ResNet-9 (\textit{tiny\_resnet}), and torchvision's
        MobileNet-V3-Small with optional ImageNet transfer
        learning.
    \item Image-dataset upload: a zip of \texttt{<class>/<file>}
        folders, extracted into the canonical \texttt{ImageFolder}
        layout that PyTorch's data loaders expect.
    \item Live training visualisation reused from the tabular path
        (the \texttt{training\_epoch} socket event fans out into
        train/val loss and val-accuracy series).
    \item Single-image inference (the ``Try it'' panel)
        immediately after training completes.
\end{itemize}

What's deliberately out:

\begin{itemize}
    \item Freeform layer-builder DAG --- too easy to OOM on stage.
    \item LLM fine-tuning --- won't fit inside 6\,GB.
    \item Object detection / segmentation --- different label
        format, different demo.
    \item ONNX export --- a one-line follow-up if needed, not in
        scope.
    \item Multi-GPU / distributed training.
\end{itemize}

Drawing the line clearly was as important as anything I built. A
no-code deep-learning interface that lets the user wire up an
architecture that doesn't fit is, in practice, a debugger they
didn't ask for.

\subsection{Why a separate microservice}

Two reasons:

\begin{enumerate}
    \item \texttt{torch} plus \texttt{torchvision} add about
        $2$\,GB of dependencies that the existing H2O-based
        \texttt{ml-training-service} doesn't need. Co-locating them
        would have doubled the size of an image that rebuilds on
        every CI run.
    \item GPU access is per-container in Docker. The DL service
        declares a \texttt{deploy.resources.reservations.devices}
        block; the H2O service runs on CPU. Mixing them in one
        container would either deny GPU to the DL workload or
        force the H2O workload onto a GPU it can't use.
\end{enumerate}

The service has its own Celery queue (\texttt{dl\_training}) so a
20-epoch CIFAR run cannot starve the fast tabular queue.

\section{Implementation}

\subsection{The architecture catalogue}

\begin{itemize}
    \item \textbf{LeNet} --- roughly $60$\,k parameters. Smoke-test
        scale; trains in seconds on Fashion-MNIST-class data.
    \item \textbf{Tiny ResNet (ResNet-9)} --- roughly $2.7$\,M
        parameters. Two residual blocks, channels
        $64\to128\to256$. Hits around $90\%$ accuracy on CIFAR-10
        in five epochs at batch~64, input $64$\,px.
    \item \textbf{MobileNet-V3-Small} --- roughly $2.5$\,M
        parameters. Loaded from torchvision with optional
        ImageNet weights. The transfer-learning toggle is the
        single biggest UX win for the demo --- a fresh user's
        200-image custom dataset can hit 80\%+ accuracy in two
        minutes, where random-init would take twenty.
\end{itemize}

\subsection{The VRAM guard}

The most subtle piece of Sprint~4 is the static memory estimator:

\begin{lstlisting}[language=Python, caption=Static VRAM estimate (\texttt{dl-training-service/app/services/vram\_guard.py})]
def estimate(arch: str, input_size: int, batch_size: int) -> Estimate:
    weights = params_mb(arch)
    optimizer = 3.0 * weights              # weights + grads + 2 Adam moments
    activations = (
        _ACTIVATION_PER_PIXEL_MB[arch]
        * (input_size * input_size)
        * batch_size
    )
    total = weights + optimizer + activations + _CUDA_RUNTIME_MB
    return Estimate(weights_mb=weights, optimizer_mb=optimizer,
                    activations_mb=activations,
                    runtime_mb=_CUDA_RUNTIME_MB, total_mb=total)
\end{lstlisting}

The route refuses any \texttt{/dl/train} request whose estimate
exceeds the user's tier-resolved budget minus a 1\,GB headroom for
the CUDA runtime and allocator fragmentation. This is the single
most important demo-safety net in the system: a wrong batch-size
press cannot OOM the GPU on stage.

The constants in \texttt{\_ACTIVATION\_PER\_PIXEL\_MB} are
conservative empirical estimates. The unit-test suite enforces that
every \texttt{(arch, input\_size, batch\_size)} triple in the
documented catalogue fits on a 6\,GB card after headroom, but it
does not pin the absolute numbers --- those drift slightly between
GPU generations and need recalibration on first deploy.

\subsection{Tier-aware ceilings}

Sprint~4 also added two new per-tier columns to the
\texttt{subscriptions} table: \texttt{max\_dl\_epochs} and
\texttt{max\_dl\_batch\_size}. The defaults:

\begin{itemize}
    \item \textbf{Free} --- 5 epochs, batch 32.
    \item \textbf{Solo} --- 20 epochs, batch 64.
    \item \textbf{Company} --- 50 epochs, batch 64.
\end{itemize}

The defence-in-depth chain runs:

\begin{enumerate}
    \item The frontend's \texttt{DLTrainNode} reads the user's
        \texttt{limits} from \texttt{/users/me} and clamps the
        slider \texttt{max} accordingly --- a free-tier user
        literally cannot select 20 epochs from the UI.
    \item The route enforces the tier ceiling regardless: if a
        custom client posts 50 epochs to \texttt{/dl/train} on a
        free-tier token, the route returns \texttt{400
        epochs\_over\_tier\_limit}.
    \item The \texttt{vram\_guard} then checks the static memory
        estimate against the tier's VRAM budget.
    \item The service-wide \texttt{HARD\_MAX\_EPOCHS} and
        \texttt{HARD\_MAX\_BATCH\_SIZE} caps in \texttt{config.py}
        sit as an absolute ceiling that no tier override can lift.
\end{enumerate}

Four layers may sound paranoid, but each one catches a different
class of mistake --- accidental, malicious, oversized, or
pathological --- and none of them is expensive.

\subsection{Frontend integration}

\figph{Pipeline editor screenshot --- DL mode}{Pipeline editor in deep-learning mode}

The canvas gained a third toggle button (``Deep Learning''), three
new node components (\texttt{ImageDatasetNode}, \texttt{CNNArchNode},
\texttt{DLTrainNode}), and a fourth preset (\texttt{dlStarter}).
The three-family lock logic was extended to lock all three
families against each other once any node of a given family lands
on the canvas.

The \texttt{ImageDatasetNode} body shows class chips, total-image
count, and a deterministically seeded thumbnail strip (three images
per class, pulled from \texttt{/datasets/<id>/image-preview} as
base64-encoded JPEGs). The \texttt{CNNArchNode} surfaces a
\texttt{pretrained} toggle that disables itself when the chosen
architecture has no canonical ImageNet checkpoint --- the user sees
\textit{why} the option is unavailable, rather than having a
control silently disappear, which is the kind of small UX detail
that tends to define whether people trust the interface.

\figph{ImageDatasetNode close-up}{ImageDatasetNode --- class chips, image count, thumbnail strip}

\subsection{Live progress and inference}

The training Celery task emits a \texttt{training\_epoch} socket
event per epoch with \texttt{\{epoch, train\_loss, val\_loss,
val\_acc\}}. The \texttt{LiveTrainingChart} component fans that
single event into three Plotly series, all sharing the same
epoch-indexed X-axis.

When training completes, a \texttt{DLPredictPanel} component
renders inline, immediately under the success alert. The user drops
an image, the panel POSTs it to \texttt{/dl/predict/<version\_id>},
and the top-five softmax probabilities render as ranked horizontal
bars.

\figph{DLPredictPanel screenshot}{Inline DL prediction panel --- top-five softmax probabilities}

\subsection{Demo path end-to-end}

For the live demo, the path is:

\begin{enumerate}
    \item \texttt{python scripts/seed\_demo\_image\_dataset.py
        /tmp/demo.zip} --- generates a 30-image, 3-class synthetic
        dataset (red circles, green squares, blue triangles).
    \item Upload the zip from the Data page.
    \item Drop the \texttt{dlStarter} preset onto a fresh pipeline.
    \item Click Run. Roughly 30 seconds later, the predict panel
        renders.
    \item Drop an image --- one of the synthetic ones, or a real
        one from a phone --- and inference returns in roughly
        $5$\,ms on CPU.
\end{enumerate}

A separate \texttt{scripts/dl\_smoke.sh} script exercises the
entire chain non-interactively: gateway reachability, GPU
visibility, the oversize-refusal path, the demo run, and the
polling loop. Having that script saved me on more than one
occasion when something seemingly unrelated had broken the DL
path.

\section{Unit Tests}

The \texttt{dl-training-service} ships with 22 unit tests:

\begin{itemize}
    \item Architecture forward-pass shapes --- one test per
        architecture on a known input tensor.
    \item VRAM-guard envelope --- parametrised over every
        documented \texttt{(arch, input\_size, batch\_size)}
        triple, asserts the estimate is under the 6\,GB-minus-1\,GB
        budget.
    \item Oversize refusal --- a request with an obviously
        oversized batch is rejected by the route before training
        starts.
    \item End-to-end \texttt{build $\to$ save $\to$ predict} on
        a tiny LeNet against a 12-image fixture dataset, against
        a real torch instance (CPU on CI).
    \item Tier-ceiling enforcement --- a free-tier token posting
        20 epochs is rejected with the documented error code.
\end{itemize}

\section*{Conclusion}

Sprint~4 was the riskiest sprint of the project. Bolting an
entirely new workload onto a platform that had calmed down into
three sprints of mostly-tabular and RAG work could have
destabilised the whole thing. It didn't, mostly because the
conventions established in earlier sprints --- the shared-volume
pattern, the Celery-per-service split, the gateway header injection
--- absorbed the new service without needing to bend.

% ═════════════════════════════════════════════════════════════════════════════
\chapter{Sprint 5: Multi-tenancy, Project ACL, and Real-Time Collaboration}
% ═════════════════════════════════════════════════════════════════════════════

\section*{Introduction}

Up to this point the platform had been a single-user experience.
The company-tier story needed something else: multiple people on
the same pipeline at the same time, with sensible access controls
and enough real-time feedback that they wouldn't accidentally undo
each other's work.

\section{Sprint Backlog}

\begin{itemize}
    \item Per-pipeline access control with five distinct roles.
    \item Server-side and client-side role enforcement.
    \item Real-time canvas presence (cursors and node highlights).
    \item Threaded chat per pipeline.
    \item One-click Google Meet links (via Google OAuth).
    \item Per-pipeline notification fan-out (the bell in Sprint~6
        will subscribe to this).
\end{itemize}

\section{Design}

\figph{Use case diagram for Sprint 5}{Use case diagram --- collaboration roles and actions}

\figph{Sequence diagram --- real-time presence}{Sequence diagram --- cursor + node-selection events over SocketIO}

\figph{Class diagram for Sprint 5}{Class diagram --- pipeline membership, chat threads, Meet integration}

\subsection{The role hierarchy}

\begin{itemize}
    \item \textbf{Owner} --- the user who created the pipeline.
        Cannot be removed.
    \item \textbf{Project Manager} --- can invite, can edit, cannot
        remove the owner.
    \item \textbf{Data Scientist} --- can edit nodes, can run
        training, cannot manage members.
    \item \textbf{Analyst} --- can view metrics and download
        artefacts, cannot run training.
    \item \textbf{Viewer} --- read-only.
\end{itemize}

The role is enforced both server-side and client-side. The gateway
resolves the user's role for the project on every
\texttt{/pipelines/*} request and stamps \texttt{X-Project-Role};
the frontend disables UI affordances based on the same role. The
server-side enforcement is authoritative --- the client-side
disable is for UX, not security. Mixing those two up is one of the
easier mistakes to make in this area.

\figph{Manage Access dialog}{Manage Access dialog --- invite, role change, remove}

\section{Implementation}

\subsection{Real-time presence}

When a user opens a pipeline, the gateway joins them to the
\texttt{pipeline\_<id>} SocketIO room. Cursor positions and node
selections are emitted as throttled (50\,ms) events. The frontend
maintains a per-user coloured cursor overlay on the canvas. There's
a particular satisfaction in seeing another user's cursor move
across your screen --- it makes the platform feel inhabited in a
way that chat alone never quite does.

\figph{Real-time presence screenshot}{Real-time presence --- two users editing the same pipeline}

\subsection{Pipeline chat and Google Meet}

Each pipeline gets its own threaded chat, persisted to Postgres. An
``@'' mention fires a notification (which the Sprint~6 bell will
have somewhere to surface). The Google Meet integration relies on
Google's OAuth flow --- not for any data-sharing reason, but
because a Meet link generated server-side without OAuth wouldn't be
tied to the inviting user's calendar, which made it useless for the
actual workflow.

\subsection{Why the gateway resolves the role}

The role-resolution lookup happens inside the gateway rather than
each downstream service. Two reasons:

\begin{enumerate}
    \item The gateway already maintains the request-scoped user
        identity. Adding role resolution here keeps the
        zero-trust pattern intact --- downstream services trust
        the stamped \texttt{X-Project-Role} the same way they
        trust \texttt{X-User-Id}.
    \item It centralises the cache. A short-lived Redis-backed
        cache on \texttt{(user\_id, pipeline\_id) $\to$ role}
        cuts the auth-service lookup count to a small fraction
        of the request count, without spreading caching logic
        across four services.
\end{enumerate}

\section{Unit Tests}

\begin{itemize}
    \item Role-matrix table-test: for each
        \texttt{(role, action)} pair, the expected
        permit/deny outcome.
    \item Role demotion mid-session: a Project Manager has their
        role downgraded to Viewer while connected; the next
        canvas-edit event is rejected.
    \item Owner immutability: an attempt to remove the owner
        returns the documented error code.
    \item SocketIO presence: two test clients join the same
        pipeline room and assert that each sees the other's
        cursor events.
\end{itemize}

\section*{Conclusion}

The collaboration layer turned the platform from a personal tool
into something a team could share. The interesting part of this
sprint was how little of the existing architecture had to change to
support it --- the gateway already injected user identity, the
SocketIO bus was already in place for live training, and the role
check fit naturally into the existing header convention.

% ═════════════════════════════════════════════════════════════════════════════
\chapter{Sprint 6: Dashboard, i18n, Accessibility, and Admin Console}
% ═════════════════════════════════════════════════════════════════════════════

\section*{Introduction}

Every project has the sprint where the punch list of small things
finally gets crossed off. This was that sprint. The dashboard, the
billing UI, internationalisation, the high-contrast theme, the
notification bell, the admin operations dashboard, user
impersonation, and undo/redo on the canvas --- each modest on its
own, but together they moved the platform from ``demonstrably
works'' to ``demonstrably operable''.

\section{Sprint Backlog}

\begin{itemize}
    \item User-facing dashboard (datasets, pipelines, model
        versions in one view).
    \item Billing UI (free / solo / company tier; Stripe
        checkout).
    \item Internationalisation (English + French, generated
        bindings).
    \item High-contrast theme + WCAG AA audit.
    \item Persistent notification bell.
    \item Admin operations dashboard (user management, audit log,
        queue + hardware monitors, migration drift panel).
    \item User impersonation with hard five-minute TTL.
    \item Undo / redo on the canvas.
\end{itemize}

\section{Design}

\figph{Use case diagram for Sprint 6}{Use case diagram --- dashboard + admin operations}

\figph{Dashboard wireframe}{Wireframe --- user dashboard with datasets, pipelines, versions}

\figph{Class diagram for Sprint 6}{Class diagram --- notifications, audit log, impersonation token}

\section{Implementation}

\subsection{The user dashboard and billing UI}

\figph{Dashboard screenshot}{User dashboard --- datasets, pipelines, model versions, quota}

\figph{Billing page screenshot}{Billing page --- tier selection, Stripe checkout}

The dashboard is a single page with three tabs (Datasets,
Pipelines, Models), each backed by a paginated table with
server-side sorting. The right rail shows the user's quota
(datasets, model versions, RAG documents) for the current tier and
a ``Manage Plan'' link to the billing UI. Billing itself goes
through Stripe Checkout for purchases and the Stripe Customer
Portal for upgrades and cancellations --- the platform never
handles a card number directly.

\subsection{Internationalisation}

Every user-facing string moved to \texttt{i18next}, with English
and French resource bundles. The TypeScript bindings are generated
from the English bundle, so a missing French key fails the
type-check rather than rendering as a raw key string in
production. That's a small piece of automation that paid for itself
the first time I forgot to translate a label.

\figph{Language toggle in action}{Language toggle --- English / French}

\subsection{High-contrast theme and accessibility}

A high-contrast theme was added for users with visual impairments,
toggleable from the navbar. Every interactive element gained an
\texttt{aria-label}; the colour palette was checked against WCAG~AA
contrast ratios. Accessibility isn't an optional extra at the
company tier --- in some jurisdictions it's a procurement
requirement.

\subsection{Notification bell with persistence}

A bell icon in the navbar shows unread chat mentions, training-done
events, training-failed events, document-indexed events, and
meeting links. The store is persisted to \texttt{localStorage} so a
page refresh doesn't drop the unread count, which sounds obvious
until you build the first version and the count keeps resetting to
zero.

\subsection{Admin operations dashboard}

The super-admin role gained a dedicated page with five panels:

\begin{itemize}
    \item \textbf{User management} --- search, suspend, delete,
        override quotas.
    \item \textbf{Audit log viewer} --- paginated table of every
        sensitive action (login, password reset, impersonation,
        suspension), with IP and JSON detail.
    \item \textbf{Queue monitor} --- live stats from each Celery
        queue.
    \item \textbf{Hardware monitor} --- CPU, RAM, GPU utilisation
        sampled from the metrics-service every five seconds.
    \item \textbf{Migration drift panel} --- compares each
        service's latest Alembic revision against the database
        and surfaces the diff.
\end{itemize}

\figph{Admin dashboard screenshot}{Admin operations dashboard --- five-panel layout}

\subsection{User impersonation}

Super-admins can ``view as'' any user from the user-management
table. The auth-service mints a separate five-minute access token
with an \texttt{imp\_actor} claim identifying the original
super-admin. A red banner persists across every page during the
impersonation session, showing the target email and a live
countdown to expiry. Clicking \texttt{Exit} (or hitting Logout in
the navbar) restores the super-admin's original token. Two audit
events are recorded: \texttt{admin.impersonate\_start} and
\texttt{admin.impersonate\_end}.

The hard TTL is the load-bearing piece here. An impersonation
session that doesn't time out is, sooner or later, an incident
waiting to happen.

\figph{Impersonation banner}{Impersonation banner --- target email + live countdown}

\subsection{Undo / redo on the canvas}

A 50-snapshot history stack with drag-coalescing (a single drag
gesture produces one snapshot, not 60) and
\texttt{Cmd}/\texttt{Ctrl+Z} plus \texttt{Cmd}/\texttt{Ctrl+Shift+Z}
keyboard shortcuts. Two \texttt{IconButton}s sit next to the
auto-layout button for users who don't know the keyboard shortcut.
It's a small feature, but it's one of those things you don't notice
until it's missing.

\subsection{Cross-cutting changes}

Two architectural shifts landed this sprint that don't fit neatly
under any one feature:

\begin{itemize}
    \item \textbf{Auth state moved to \texttt{sessionStorage}.}
        Before this sprint, refreshing the page during
        impersonation lost the original super-admin token --- the
        Zustand auth store was in-memory only. Wrapping it in
        \texttt{persist(\ldots,\,sessionStorage)} fixed the
        impersonation case and quietly improved the general
        refresh behaviour for normal users at the same time.
    \item \textbf{ConfirmDialog with typed confirmation.}
        Destructive admin actions (delete-user,
        delete-announcement) now require retyping the target's
        email or title before the action button enables. The
        native \texttt{confirm()} dialog was too easy to misclick
        during a live demo, and once was enough.
\end{itemize}

\section{Unit Tests}

\begin{itemize}
    \item i18n key parity: a CI step diffs the English and French
        bundles and fails if a key is missing on either side.
    \item Notification reducer: state transitions for each event
        type (mention, training-done, etc.), plus persistence
        round-trip.
    \item Impersonation flow: the original token, the
        impersonation token, and the post-exit restore are
        asserted across three pytest cases.
    \item Audit-log presence: every sensitive route is exercised
        and the audit table is asserted non-empty for that
        action.
    \item Undo / redo: a sequence of canvas edits, then a fixed
        number of \texttt{Ctrl+Z} presses, then a deep-equality
        check against an earlier snapshot.
\end{itemize}

\section*{Conclusion}

The punch-list sprint never makes for thrilling reading, but it's
where the platform earned the right to be presented to a real
user. The admin console in particular changed how I thought about
the platform --- having a single place to inspect queue health,
hardware load, and migration state turned operations from a vague
worry into a tractable one.

% ═════════════════════════════════════════════════════════════════════════════
\chapter{Testing and Deployment}
% ═════════════════════════════════════════════════════════════════════════════

\section*{Introduction}

The platform's testing and deployment stories evolved in parallel
with the rest of the system, and by the end they looked nothing
like what I had at the start. This chapter covers both --- how the
tests are organised, what they actually verify, and how the whole
thing ships.

\section{Test Strategy}

\subsection{Backend Testing}

The test suite is split by service. Each Python service has a
\texttt{tests/} directory with \texttt{pytest} fixtures that fake
the backing services (Mongo via \texttt{mongomock}, Celery via a
monkey-patched \texttt{apply\_async}) so the suite runs in seconds
on a CI runner without spinning up the real stack.

Coverage is unevenly distributed, which I won't pretend otherwise
about:

\begin{itemize}
    \item \texttt{auth-service} --- the authentication path and
        the admin operations sit above $80\%$ line coverage.
    \item \texttt{ml-training-service} --- the algorithm-factory
        path is fully exercised; the live-streaming SocketIO
        path is not.
    \item \texttt{dl-training-service} --- 22 unit tests covering
        the route's contract surface, the architecture
        forward-pass shapes, the VRAM-guard envelope, and an
        end-to-end \texttt{build $\to$ save $\to$ predict} happy
        path on a real torch instance.
\end{itemize}

\figph{Backend coverage report}{Backend coverage report --- per-service line coverage}

\subsection{Frontend Testing}

The frontend uses \texttt{vitest} for unit tests but relies mostly
on TypeScript strict mode plus ESLint as the structural guarantee.
In practice more than $99\%$ of regressions in this codebase are
caught by \texttt{tsc --noEmit} before they ever reach a runner ---
typing is the cheapest test you can write.

\figph{Frontend coverage report}{Frontend coverage report --- vitest output}

\subsection{CI runs}

\figph{GitHub Actions --- backend job}{GitHub Actions output for the backend lint + test job}

\figph{GitHub Actions --- frontend job}{GitHub Actions output for the frontend lint + test job}

End-to-end, CI runs in roughly six minutes for a frontend-only
change, nine minutes for a backend change, and fourteen minutes
when both sides move at once.

\section{Deployment}

The whole stack ships as a single \texttt{docker-compose.yml} with
14 services:

\begin{itemize}
    \item 4 backing data stores (Postgres + pgvector, Mongo,
        Redis, TimescaleDB).
    \item 5 application services (gateway, auth, data, ml, dl).
    \item 3 Celery workers, one per CPU-bound service.
    \item 1 metrics collector.
    \item Ollama plus a one-shot \texttt{ollama-pull} init
        container.
    \item MailHog for development email capture.
\end{itemize}

A \texttt{make up} brings the entire platform online. First-run
duration is dominated by the model pulls: Ollama fetches the LLM
weights on first boot (around $2$\,GB), and the dl-training
service pulls the CUDA wheels on first build (around $1$\,GB).
After that, \texttt{make up} is a matter of seconds.

\figph{Deployment diagram}{Deployment diagram --- 14 containers, three named volumes}

\subsection{Volumes}

Three named Docker volumes carry the binary artefacts that survive
container restarts: \texttt{uploaded\_files} (datasets and image
zips), \texttt{model\_artifacts} (trained models), and
\texttt{ollama\_models} (LLM weights). \texttt{uploaded\_files} is
mounted into the gateway, both ingestion containers (API and
worker), the ml-training containers, and the dl-training
containers --- the shared-volume convention from Sprint~1 ended up
applied very broadly, arguably more broadly than I'd originally
planned.

\subsection{Healthchecks and dependency ordering}

Every backing service has a \texttt{healthcheck}, and every
application service \texttt{depends\_on} the relevant healthchecks
with \texttt{condition: service\_healthy}. The first time I
deployed the stack without this discipline, the auth service
started before Postgres had finished initialising and proceeded to
crash-loop on the Alembic migration. After that, the healthcheck
convention became non-negotiable for any service that touches a
database.

\subsection{Network and security topology}

\figph{Network topology}{Network topology --- gateway as the single ingress, internal Docker network for the rest}

The gateway is the only container exposed on the host. All other
services bind to the internal Docker network. The Postgres
instance is not reachable from outside the Compose stack; the only
path to it is through the auth-service or the ml-training-service,
and the only path to either of those is through the gateway.
That is the same zero-trust posture as the JWT header injection,
applied at the network layer.

\subsection{Bugs encountered during deployment}

A non-exhaustive list, in roughly the order they bit me:

\begin{itemize}
    \item \textbf{Prophet missing CmdStan.} Prophet's probabilistic
        backend is a C++ binary that has to be compiled at
        image-build time, not pip-installed. The fix was a
        \texttt{python -c "import\,cmdstanpy;\,cmdstanpy.install\_cmdstan()"}
        step in the ml-training Dockerfile.
    \item \textbf{Flask CLI app discovery inside containers.}
        Running \texttt{flask db migrate} via \texttt{docker
        compose exec} couldn't find the app factory; the fix
        was an explicit \texttt{-e FLASK\_APP=app.main} on every
        CLI invocation, baked into the Makefile so I wouldn't
        forget.
    \item \textbf{React 18 Strict Mode and rate limiting.} Strict
        Mode double-invokes \texttt{useEffect} in development;
        the rate limit was tripping. Bumped the per-user limit
        to $600$/min in development, kept the production cap
        tight.
    \item \textbf{torch CUDA wheels in CI.} The
        \texttt{dl-training-service} Dockerfile pulls CUDA-12.1
        torch from PyTorch's index; on a CI runner that would
        have been a $1$\,GB download for no benefit. CI installs
        CPU-only torch wheels separately, gated on
        \texttt{matrix.service == 'dl-training-service'}.
    \item \textbf{Ollama GPU detection in Compose.} The default
        compose deployment runs Ollama on CPU. Enabling the
        \texttt{deploy.resources.reservations.devices} block
        requires \texttt{nvidia-container-toolkit} on the host,
        which isn't the default on Ubuntu Server. Documented in
        the README and commented out in the compose file by
        default.
    \item \textbf{DL VRAM-guard calibration.} The static
        per-pixel constants were tuned against documented Turing
        memory numbers rather than measured ones. The first real
        GPU run on my 1660~Super exceeded the estimate by about
        $15\%$ on MobileNet at 224\,px batch~32; the constants
        got nudged upwards accordingly. The unit-test suite
        tracks the envelope shape, not the absolute MB, so this
        calibration drift was caught empirically rather than by
        a failing test.
\end{itemize}

\section*{Conclusion}

The stack is, by deployment-engineering standards, modest --- one
host, fourteen containers, \texttt{make up}. That simplicity is
the point. A platform whose deployment story requires Kubernetes
expertise wouldn't be a credible answer to the question this
project set out to ask.

% ═════════════════════════════════════════════════════════════════════════════
\chapter*{General Conclusion}
\addcontentsline{toc}{chapter}{General Conclusion}
% ═════════════════════════════════════════════════════════════════════════════

The end-of-studies project set out to build a privacy-first,
no-code AI platform that could stand alongside cloud-managed ML
services without making the data-sovereignty compromises those
services demand. After six sprints, the platform supports three
pipeline modes --- tabular ML with fourteen algorithms,
fully-local RAG-augmented chat, and image-classification deep
learning --- runs on a single consumer machine with $6$\,GB of
VRAM, and exposes that capability through a visual canvas designed
for users who don't write Python.

\section*{Key deliverables}

\begin{itemize}
    \item \textbf{Five Python microservices} totalling roughly
        $6{,}500$ lines of code across the api-gateway,
        auth-service, data-ingestion-service, ml-training-service,
        metrics-service, and dl-training-service. Three
        independent Celery worker processes handle the
        asynchronous workloads.
    \item \textbf{A React + TypeScript single-page application} of
        roughly $13{,}000$ lines, including the React Flow
        canvas, the live training visualisations, the admin
        operations dashboard, the i18n bundles for English and
        French, the high-contrast theme, and the impersonation
        banner.
    \item \textbf{14 docker-compose services} (5 application
        services, 3 Celery workers, 4 backing data stores,
        Ollama, \texttt{ollama-pull}, MailHog).
    \item \textbf{14 tabular ML algorithms} covering
        classification, regression, clustering, and time-series
        forecasting.
    \item \textbf{A fully-local RAG pipeline} built on
        sentence-transformers, pgvector, and Ollama serving
        Llama~3.2~3B \texttt{Q4\_K\_M}.
    \item \textbf{3 deep-learning architectures} (LeNet, custom
        ResNet-9, MobileNet-V3-Small with optional ImageNet
        transfer learning), guarded by a static VRAM estimator
        that refuses oversize requests before the GPU sees them.
    \item \textbf{A GitHub Actions CI pipeline} with parallel
        lint and test jobs across all five Python services and
        the TypeScript frontend.
    \item \textbf{An Alembic migration history} of nine
        revisions across two services, idempotent and
        replayable.
\end{itemize}

\section*{Validation results}

The platform was validated end-to-end on a representative tabular
classification problem (a $250{,}000$-row salary dataset, Random
Forest regression, $R^2$ of $0.9785$ in $4$\,min\,$12$\,s on the
demo machine) and on a representative DL classification problem
(the synthetic 30-image / 3-class dataset shipped as
\texttt{seed\_demo\_image\_dataset.py}, tiny\_resnet at $64$\,px
batch~32, $5$ epochs, terminal validation accuracy of $100\%$ in
roughly $30$\,s on a 1660 Super).

\figph{Final results dashboard}{Final results dashboard --- tabular + DL validation runs}

\section*{What worked, what I'd change}

\textbf{What worked.} The single design decision that paid off
most was the gateway-only header-injection rule from Sprint~1.
Every downstream service can trust \texttt{X-User-Id} without
inspection, which removed an entire category of cross-service auth
bugs that a more permissive setup would have invited. The
shared-volume convention was the second --- it saved me from
writing a streaming-upload protocol, and it generalised cleanly to
RAG documents, image zips, and trained-model artefacts. The
Celery-per-service split was the third --- a long XGBoost run
genuinely cannot block an authentication request, and a 20-epoch
DL run cannot starve the fast tabular queue.

\textbf{What I'd change.} Flask over FastAPI: the absence of typed
routes cost me hours over the project's life, especially in
\texttt{ml-training-service} where the request bodies are large
and nested. Three algorithm registries instead of one: the tabular
registry, the architecture registry, and the algorithm-display
catalogue all encode mostly the same information in slightly
different shapes; a single config-driven catalogue would have been
easier to maintain. A formal schema for socket events: the
\texttt{/training} namespace grew organically and is now mildly
inconsistent (some events use \texttt{step}, some use
\texttt{epoch}, some use both). Lesson learned the slow way.

\section*{Limitations and future work}

The platform is single-host by construction --- the shared-volume
convention assumes a single filesystem. Scaling out the workers
would need either a network filesystem or an S3-compatible object
store. Per-user quota overrides require a JWT refresh to take
effect, because the tier rides in the JWT claim. The DL VRAM
guard is calibrated against the demo machine; a truly portable
estimator would instrument a one-time calibration run on first
boot.

Future work: ONNX export for DL versions (the dl-training service
already saves a clean state\_dict --- this is a one-file change
away), object detection via a pre-built YOLOv8 model for inference
only, per-user quota override invalidation via a Redis-backed
override cache the gateway consults, and a typed schema for the
\texttt{LiveTrainingChart} socket events.

\section*{Closing remark}

What I take away from this project is a specific belief: the
common framing that ``cloud AI is the only way you can scale ML''
is not exactly wrong, but it understates how much can be done
locally. A laptop with $6$\,GB of VRAM, a competent CPU, and
$16$\,GB of RAM is enough to host an AutoML platform, a small
generative-AI workspace, and a deep-learning training loop --- all
at once, and without any of the data ever leaving the machine. For
the regulated industries this project was framed against, that's
not a marginal improvement. It's the difference between machine
learning being usable at all and not.

% ═════════════════════════════════════════════════════════════════════════════
\begin{thebibliography}{99}

\bibitem{scikit-learn}
F. Pedregosa et al.,
\textit{Scikit-learn: Machine Learning in Python},
Journal of Machine Learning Research, vol.~12, pp.~2825--2830, 2011.

\bibitem{xgboost}
T. Chen and C. Guestrin,
\textit{XGBoost: A Scalable Tree Boosting System},
Proceedings of the 22nd ACM SIGKDD International Conference on
Knowledge Discovery and Data Mining, pp.~785--794, 2016.

\bibitem{shap}
S. M. Lundberg and S.-I. Lee,
\textit{A Unified Approach to Interpreting Model Predictions},
Advances in Neural Information Processing Systems 30 (NeurIPS 2017).

\bibitem{prophet}
S. J. Taylor and B. Letham,
\textit{Forecasting at Scale}, The American Statistician, 2018.

\bibitem{minilm}
W. Wang et al.,
\textit{MiniLM: Deep Self-Attention Distillation for Task-Agnostic
Compression of Pre-Trained Transformers}, NeurIPS 2020.

\bibitem{rag}
P. Lewis et al.,
\textit{Retrieval-Augmented Generation for Knowledge-Intensive NLP
Tasks}, NeurIPS 2020.

\bibitem{llama32}
Meta AI,
\textit{Llama 3.2: Open foundation and instruction models},
\url{https://ai.meta.com/llama/}, 2024.

\bibitem{mobilenetv3}
A. Howard et al.,
\textit{Searching for MobileNetV3},
Proceedings of the IEEE/CVF International Conference on Computer
Vision (ICCV), 2019.

\bibitem{resnet}
K. He, X. Zhang, S. Ren, and J. Sun,
\textit{Deep Residual Learning for Image Recognition},
IEEE Conference on Computer Vision and Pattern Recognition (CVPR),
2016.

\bibitem{lenet}
Y. LeCun, L. Bottou, Y. Bengio, and P. Haffner,
\textit{Gradient-Based Learning Applied to Document Recognition},
Proceedings of the IEEE, vol.~86, no.~11, pp.~2278--2324, 1998.

\bibitem{pytorch}
A. Paszke et al.,
\textit{PyTorch: An Imperative Style, High-Performance Deep
Learning Library}, NeurIPS 2019.

\bibitem{ollama}
Ollama,
\textit{Get up and running with large language models locally},
\url{https://ollama.com/}, 2024.

\bibitem{pgvector}
A. Kane,
\textit{pgvector: Open-source vector similarity search for
Postgres},
\url{https://github.com/pgvector/pgvector}, 2024.

\bibitem{reactflow}
xyflow GmbH,
\textit{React Flow Documentation},
\url{https://reactflow.dev/}, 2024.

\bibitem{celery}
Celery Project,
\textit{Celery: Distributed Task Queue},
\url{https://docs.celeryq.dev/}, 2024.

\end{thebibliography}

\end{document}
